\documentclass[11pt]{article}
\usepackage[margin=1in]{geometry}
\usepackage[T1]{fontenc}
\usepackage{lmodern,booktabs,amsmath,graphicx}
\usepackage{microtype}
\usepackage[hidelinks]{hyperref}
\usepackage{xurl}
\graphicspath{{simulation_output/perception_precision_20261002/}{./}}
\setlength{\parindent}{0pt}
\setlength{\parskip}{6pt}
\setlength{\emergencystretch}{2em}
\title{Improving Native Perception Inputs to the NRMM Estimator}
\author{Pan Dynamics Collision Avoidance Research}
\date{October 2, 2026}
\begin{document}
\maketitle

\section{Finding and implementation}
On the reused native CARLA development recording, preserving target identity
and retaining visible near/lower vehicle returns reduced target-estimator
position RMSE from 1.24063 to 0.28514 m, and inertial velocity RMSE from
3.67216 to 0.74530 m/s. With a 0.3-second camera blackout, the selected
configuration achieved 0.33721 m and 0.89059 m/s. Both selected runs passed
the original smoke thresholds. The estimator equations, gains, initial-state
rule, approximate target dimensions and time step were unchanged.

The production change is an explicit, optional persistent image-box selector
in \path{perception/yolo_lidar_target_position.py}. The first acquisition
retains the previous confidence policy. Subsequent detections must overlap
the last accepted box; higher-confidence parked cars cannot simply replace
it. A missing or expired match emits no detection. Expiration never silently
starts another identity. Exactly one camera is required. The existing
confidence-only default remains available and unchanged.

For this native ego-CG-referenced input dataset, the selected invocation adds:
\begin{verbatim}
--camera-id front --causal --target-association persistent
--min-depth 0.1 --min-object-height 0.0
\end{verbatim}
The association minimum IoU is 0.1, maximum unmatched interval is 1.0 s,
and cluster radius remains 1.2 m. Geometry still uses a declared 5.0-by-1.9 m
sedan rectangle, a 1-degree heading grid and three causal current-frame
heading iterations. These are experiment-specific settings, not newly
validated universal defaults.

\section{Controlled comparison and scope}
The input is the preceding native CARLA recording in Town10HD\_Opt.
It contains 300 frames at 50 Hz from a front camera, LiDAR, GNSS and IMU,
with separate truth.
Its original capture and sensor configuration are documented in
\path{report/CARLA_FULL_PERCEPTION_20261002.tex}. This study did not connect
to, tick, modify, start or stop any CARLA server. It therefore consumed no
shared simulation scheduling slot.

Actual YOLO26x detections and measured LiDAR clusters were cached for paired
ablations. The cache reproduced all 300 original confidence selections and
availability decisions; maximum position difference was
$8.85\times10^{-15}$ m. An initial bitwise geometry comparison failed because
of floating-point differences, and the numerical difference is retained in
the audit. Detector weights, images, native ego inputs, dimensions and
estimator design were held fixed. Each candidate's predictions were written
and hashed before the separate scorer read truth. No truth heading, box,
position, velocity, output-offset correction or fitted vehicle size entered
prediction.

The hypotheses were explored adaptively on a recording already inspected in
the prior smoke task. This is a development experiment, not an independent
route, seed or test population. The blackout uses the same recorded drive:
only RGB frames 150--164 are replaced by a black image. Cached ablations
reuse the preceding actual detector's empty outputs there; the final selected
baseline and blackout rerun actual YOLO and point-cloud processing from the
images. Scoring retains every published estimator timestamp, including
unavailable perception frames.

\section{Ablation results}
The table uses the same $t\geq2$ s interval and original thresholds:
target position RMSE $\leq0.5$ m, target velocity RMSE $\leq1.5$ m/s,
raw position RMSE $\leq0.5$ m, raw P95 $\leq1$ m, and availability outside
injected blackout $\geq95\%$. All four original ego accuracy checks remain
unchanged. ``Lower'' means a 0 m object-height floor relative to the ego CG;
``depth'' means a 0.1 m optical-depth floor; ``tight'' means a 0.3 m
connected-component radius; ``locked'' disallows geometric heading changes
around the causal motion hint. All rows after the first two use persistence.

\begin{center}
\small
\begin{tabular}{lrrrl}
\toprule
Configuration & Raw pos. (m) & Est. pos. (m) & Est. vel. (m/s) & Gates \\
\midrule
Confidence, original & 1.246 & 1.241 & 3.672 & Fail \\
Largest image box & 0.395 & 0.467 & 1.441 & Pass$^*$ \\
Persistent only & 0.401 & 0.477 & 1.307 & Pass \\
Persistent, blackout & 0.402 & 0.529 & 1.501 & Fail \\
Locked heading & 0.405 & 0.522 & 1.485 & Fail \\
Lower returns & 0.385 & 0.411 & 0.728 & Pass \\
Lower + locked & 0.371 & 0.418 & 0.815 & Pass \\
Depth & 0.309 & 0.417 & 1.281 & Pass \\
Depth, blackout & 0.309 & 0.497 & 1.549 & Fail \\
Tight clustering & 0.376 & 0.491 & 1.311 & Pass$^*$ \\
Depth + tight & 0.277 & 0.433 & 1.364 & Pass$^*$ \\
Depth + tight, blackout & 0.273 & 0.549 & 1.747 & Fail \\
Selected: lower + depth & 0.251 & 0.285 & 0.745 & Pass \\
Selected, blackout & 0.250 & 0.337 & 0.891 & Pass \\
\bottomrule
\end{tabular}
\end{center}
$^*$Post-transient gates do not assess startup quality. Largest-box selection
excludes the intended target center on 31 valid frames and has all-frame raw
RMSE 1.591 m. Tight clustering loses frames 3--8. Depth plus tight clustering
has all-frame raw/estimator position RMSE 0.531/0.677 m, worse than the
selected configuration. These alternatives are not recommended solely because
they pass the post-transient baseline thresholds.

Selected baseline all-frame raw position RMSE is 0.22951 m and estimator
position RMSE is 0.26782 m. Blackout values are 0.22757 m and 0.30550 m.
All-frame target velocity RMSE is 1.01625 m/s for the selected baseline and
1.09143 m/s for blackout; initialization still
uses a co-moving target hypothesis. Baseline has 300/300 valid measurements;
blackout has 285/285 outside the injected interval. All 600 estimator
publications are finite. Detection returns at frame 165 immediately after
the blackout. Maximum estimator position error during the blackout is
0.78896 m, which is not a pointwise 0.5 m guarantee.

\section{Why the changes help}
\textbf{Identity continuity is the primary improvement.}
The original selector ranks each frame independently by confidence but always
emits \texttt{targetIdx=1}. In the diagnostic scoring projection, 34 original
valid boxes excluded the intended target center. Persistent selection reduces
this count to zero on this recording. This is a center-inclusion diagnostic,
not complete identity ground truth. All tested IoU gates from 0.05 to 0.4
produce the same baseline and blackout selections; the minimum accepted IoU
is 0.91676. The benefit is not a finely tuned IoU threshold on this fixture.

\textbf{The optical-depth floor clips the near vehicle face.}
At close range, selected point minimum ego $x$ repeatedly equals about
2.55 m: the 1.55 m camera extrinsic plus the original 1 m optical-depth
floor. Lowering that floor to the projector's existing 0.1 m positive-depth
limit changes mean longitudinal error from 0.27175 to 0.03363 m with all
other persistent settings fixed. This removes an input-selection bias;
subtracting a truth-fitted constant from outputs is unnecessary.

\textbf{Height selection affects temporal quality as well as raw RMSE.}
With an ego-CG reference, a 0.3 m height floor removes lower vehicle returns.
Keeping returns down to CG height reduces estimator velocity RMSE from
1.30703 to 0.72824 m/s even though raw position RMSE changes only from
0.40123 to 0.38495 m. Coupling this selection with the corrected optical-depth
floor gives the selected result. The numerical improvement is established
for this mounting/road geometry; a general ground filter still requires
road-plane estimation and a documented sensor reference point.

\textbf{Tighter clustering has a startup tradeoff.}
A 1.2 m connection radius can merge side or background points at close range;
a representative cluster extends laterally to 2.68 m. A 0.3 m radius removes
the approximately 1 m post-transient lateral peak, but fragments sparse
startup geometry, creates six unavailable measurements and worsens startup
and blackout estimation. Range-dependent clustering and visibility-aware
support tests should be studied before changing its default. Locking heading
alone also fails to improve the complete estimator consistently.

\section{Validation and remaining work}
Fourteen Python behavior tests pass: the five existing causality checks plus
nine association tests covering confidence distractors, short gaps,
expiration without identity replacement, initial missing data, camera
identity, timestamp/gate validation, explicit opt-in and the single-camera
restriction. Fresh actual detector replays agree with cached prototypes to
floating-point precision. The selected baseline and blackout have identical
first 150 perception outputs and estimator states, excluding runtime timing.
All ego states exactly match the original passing control; estimator source
and configuration hashes are unchanged.

The selected raw lateral mean error remains approximately 0.157 m, and raw
maximum error remains approximately 0.949 m. Known-size rectangle anchoring,
partial visibility and residual background points remain concrete candidates
for further investigation. Only 64/300 baseline and 59/285 blackout valid
raw measurements satisfy the configured 0.12 m error allowance. The runtime
reports available bounds on every selected frame, but that flag does not
verify physical sensor-error assumptions. Native GNSS velocity error and
out-of-domain physical target speed from the previous study are unchanged.
No domain-wide certificate or real-time claim is made.

Next validation should freeze the selected settings and use independent
ranges, turns, object crossings, initial target choices and occlusions.
Persistent image overlap alone cannot resolve long occlusions or ambiguous
crossings, and its initial confidence choice may select the wrong car in a
different scene. A multi-object association system using motion and appearance,
visibility-aware vehicle-center fitting, and a calibrated ground/reference
model are further algorithm work; none is claimed implemented here.

The tracked diagnostics are in
\path{report/PERCEPTION_PRECISION_EXPLORATION_20261002/}. Local executable
studies, cached clusters, all candidate outputs, logs and the accuracy plot
remain in \path{simulation_output/perception_precision_20261002/} and are
retained as identified export copies in the external archive bundle
\path{05_Project_B_Collision_Avoidance/2026-10-02_Perception_Precision_Exploration/}.
No raw dataset, CARLA tooling, dependency tree or generated binary is committed.

\IfFileExists{simulation_output/perception_precision_20261002/accuracy_trace.pdf}{
\clearpage
\section{Error traces on the reused recording}
\includegraphics[width=\linewidth]{accuracy_trace.pdf}
}{\IfFileExists{accuracy_trace.pdf}{
\clearpage
\section{Error traces on the reused recording}
\includegraphics[width=\linewidth]{accuracy_trace.pdf}
}{}}
\end{document}
